% ---- PREAMBLE ----
% layout
\documentclass{article}
\title{Feasibility of SSH-Free Batch Job Management in Cylc}
\author{Hilary Oliver}

\usepackage[margin=3cm, headheight=1cm]{geometry}

% Font stuff.
% \renewcommand{\familydefault}{\sfdefault}  % sanz-serif
\usepackage{parskip}  % regular paragraph separation

% Imports.
\usepackage[usenames]{color}
\usepackage{graphicx}
\usepackage{listings}
  \usepackage{courier}
\usepackage{hyperref}
\usepackage{textcomp}
\usepackage{dirtree}

% Hyperlinks.
\definecolor{links}{rgb}{0.1,0.1,0.6}
\hypersetup{colorlinks=true, linkcolor=links, urlcolor=links}
\urlstyle{same}

% Headers / Footers.
\usepackage{fancyhdr}
%\lhead{foo}
%\rhead{
%  \includegraphics[width=0.15\textwidth]{resources/png/rose-logo}
%  \includegraphics[width=0.15\textwidth]{resources/tex/cylc-logo}}
\pagestyle{fancy}

% Code listings: default style.
\definecolor{keywords}{rgb}{0.8,0.4,0.0}
\definecolor{comments}{rgb}{1.0,0.3,0.5}
\definecolor{identifiers}{rgb}{0.3,0.4,0.5}
\definecolor{strings}{rgb}{0.2,0.5,0.3}
\definecolor{basic}{rgb}{0.2,0.3,0.4}
\definecolor{command}{rgb}{0.0,0.2,0.1}
\definecolor{transcr}{rgb}{0.0,0.2,0.4}
\newcommand\mysmall{\fontsize{8}{9.2}\selectfont}

% bold for courier font:
\renewcommand{\ttdefault}{pcr}

\lstset{
basicstyle=\color{basic}\mysmall\ttfamily\bfseries,
identifierstyle=\color{identifiers},
keywordstyle=\color{keywords},
commentstyle=\color{comments},
stringstyle=\color{strings},
showstringspaces=false,
upquote=true,
}

% Code listings: suite.rc language support.
\definecolor{level1}{rgb}{0.0,0.2,0.6}
\definecolor{level2}{rgb}{0.0,0.3,0.7}
\definecolor{level3}{rgb}{0.0,0.4,0.8}
\definecolor{jinja2}{rgb}{0.7,0.5,0.3}
\lstdefinelanguage{suiterc}
{
string=[b]{"},
sensitive=true,
comment=[l]{\#},
morecomment=[s][\color{level1}]{[}{]},
morecomment=[s][\color{level2}]{[[}{]]},
morecomment=[s][\color{level3}]{[[[}{]]]},
morecomment=[s][\color{jinja2}]{\{\%}{\%\}},
morecomment=[s][\color{jinja2}]{\{\{}{\}\}},
morecomment=[s][\color{jinja2}]{\{\#}{\#\}},
}

\lstdefinelanguage{jobhosts}
{
string=[b]{'},
sensitive=true,
comment=[l]{\#},
keywords={SSH, rm, mkdir, bash, scp, rsync, hobo@otherhost, vagrant@localhost},
}

\lstdefinelanguage{transcript}
{
string=[b]{"},
comment=[l][\color{comments}]{\#},
alsoletter={$},
morekeywords={shell$},
}

\definecolor{note}{rgb}{0.6,0.6,0.6}
\newcommand{\note}[1]{{\color{note}\textbf{Note:} \textit{#1}}}

\definecolor{terminology}{rgb}{0.7,0.5,0.2}
\newcommand{\terminology}[1]{{\color{terminology} \textit{#1}}}

\definecolor{todo}{rgb}{1.0,1.0,0.0}
\newcommand{\TODO}[1]{{\colorbox{todo}{\textbf{TODO:} \textit{#1}}}}

\usepackage{framed}
\definecolor{shadecolor}{rgb}{0.95,0.95,0.95}

%\usepackage{draftwatermark}
%\SetWatermarkText{Draft 15/03/17}
%\SetWatermarkScale{5}
%\SetWatermarkColor[rgb]{0.95,0.90,0.90}

\begin{document}
\maketitle

[v1.1 - minor introductory text changes]

\section{Introduction}

Managing tasks in a workflow requires more than just job execution: cylc
performs additional actions with \lstinline=rsync= for file transfer, and
direct execution of \lstinline=cylc= sub-commands over non-interactive
SSH.\footnote{Cylc used to run bare shell expressions over SSH, which required
a bash shell and made whitelisting difficult.}

Some sites may want to restrict access to job hosts by whitelisting SSH
connections to allow only \lstinline=rsync= for file transfer, and providing a
local batch system that sees the job hosts.\footnote{A malicious script could
be \lstinline=rsync='d and executed from a batch job, but batch jobs are
considered easier to audit.}

This document describes cylc remote job management requirements, and considers
the feasibility of doing it without SSH when there is a local batch system that
sees the job hosts.

\section{Executive Summary}

{\em Section~\ref{Implementation} shows that SSH-free job management with a
local batch system is feasible} with some changes to cylc's job submit, poll,
and kill mechanisms; and some additional use of \lstinline=rsync=.

This needs to be an additional mode of operation, as cylc still needs
SSH for remote background jobs, and when a batch system is not available on
the suite host.

{\bf Estimated expert developer effort: approximately 1-2 weeks.}

PBS file staging was originally suggested as the basis for this development.
However it provides no advantage over \lstinline=rsync= in this context, and it
is less general. Some of our requirements are also suite- or host-specific, and
PBS staging is job-specific; and depending on the underlying file transfer
mechanism, staging may not be able to work with our standard suite directory
structure.

\section{SSH-based Job Management}

Cylc does not have persistent agent processes running on job hosts to act on
instructions received over the network\footnote{This would be bad use of HPC
resources, with complex implementation, administration, and security
requirements.} so instead we execute job management commands directly on job
hosts over SSH. Reasons for this include:
\begin{itemize}
  \item this works equally for batch system and background jobs
  \item SSH is {\em required} for background jobs and if the batch system is
    not available on the suite host
  \item {\em querying the batch system alone is not sufficient for full job
    polling functionality} because jobs can complete (and then be forgotten by
    the batch system) while the network, suite host, or suite server program is
    down (e.g.\ between suite shutdown and restart)
    \begin{itemize}
      \item to handle this we get the automatic job wrapper code to write
        job messages and exit status to {\em job status files} that are
        interrogated by suite server programs during job polling operations
      \item job status files reside on the job host, so the interrogation
        is done over SSH
    \end{itemize}
  \item job status files also hold batch system name and job ID; this is
    written by the job submit command, and read by job poll and kill commands
    (all over SSH)
\end{itemize}

The question now is, can we get everything we need including full job polling
functionality, with a local batch system and without SSH?

(As an aside, note that we can already use a local batch system if the suite
and job hosts have a shared filesystem, because that allows cylc to treat task
jobs as entirely local.\footnote{The job ID must also be valid to query and
kill the job via the local batch system. This is not the case for Slurm, unless
the \lstinline=--cluster= option is explicitly used in job query and kill
commands, otherwise the job ID is not recognized by the local Slurm instance.})

\section{A Concrete Example}

The following suite, registered as \lstinline=suitex=, is used to illustrate
our current SSH-based remote job management. It submits two jobs to a remote,
and a local task views a remote job log then polls and kills the remote jobs.
Each action is described, along with comments on doing it without SSH. 

\lstset{language=suiterc}
\begin{lstlisting}
# suite.rc
[scheduling]
   [[dependencies]]
          graph = "delayer => master & REMOTES"
[runtime]
   [[REMOTES]]
      script = "sleep 30"
       [[[remote]]]
           host = wizard
           owner = hobo
   [[remote-a, remote-b]]
       inherit = REMOTES
   [[delayer]]
      script = "sleep 10"
   [[master]]
       script = """
 sleep 5
 cylc cat-log -m c -f o $CYLC_SUITE_NAME remote-a.1
 sleep 2
 cylc poll $CYLC_SUITE_NAME REMOTES.1
 sleep 2
 cylc kill $CYLC_SUITE_NAME REMOTES.1
 sleep 2
 cylc remove $CYLC_SUITE_NAME REMOTES.1"""
\end{lstlisting}

The {\em delayer} task just separates suite start-up from remote job
submission, for clarity when watching the job host (e.g.\ with
\lstinline=watch -n 1 find ~/cylc-run/suitex=).

Global config specifies the path to the remote cylc executable, and tells cylc
not to use a remote login shell, and to retrieve job logs:
\begin{lstlisting}
# global.rc
[hosts]
   [[wizard]]
       cylc executable = /opt/bin/cylc
       retrieve job logs = True
       use login shell = False
\end{lstlisting}

On running the suite, remote job host actions were captured in the transcripts
below by wrapping the \lstinline=ssh=, \lstinline=scp=, and \lstinline=rsync=
executables in scripts that log their command lines before taking action.

\renewcommand*\DTstylecomment{\normalfont\ttfamily\color{comments}}
\renewcommand*\DTstyle{\bf\ttfamily\textcolor{identifiers}}

\subsection{create suite run directory and install source files}

Done by \lstinline=rose suite-run= before suite start-up
(the command will be migrated to cylc soon though).

\begin{itemize}
  \item with \lstinline=--new= it invokes bash over SSH and a raw shell
    expression, to delete previous-run files
  \item it invokes itself over over SSH to create top level suite directories
    and install source files
    \begin{itemize}
      \item skips installation if server UUID file is found on the job host
        (indicates a shared filesystem)
    \end{itemize}
  \item uses \lstinline=rsync= for suite source file installation
  \item (note the same directory structure is used on suite and job hosts, for
    consistency and simplicity, and because the suite host can also be a job host)
\end{itemize}

{\color{blue}{\bf SSH-free?}
\begin{itemize}
  \item directory creation and source file installation can be done with \lstinline=rsync=
  \item clean-out can too, with \lstinline=rsync= delete options and an empty
    source directory on the suite host?
  \item detect shared filesystems by attempting to \lstinline=rsync= the UUID
    file back from the job host?
\end{itemize} 
}

\lstset{breaklines=true}
\lstset{language=jobhosts}

\vspace{5mm}
\begin{lstlisting}
# rose suite-run --new only: initial clean-out
ssh -oBatchMode=yes -oConnectTimeout=10 hobo@wizard bash -l -O extglob -c 'cd; echo '"'"'673d7a0d-7816-42a4-8132-4b1ab394349c'"'"'; ls -d -r cylc-run/suitex/work cylc-run/suitex/share/cycle cylc-run/suitex/share cylc-run/suitex; rm -fr cylc-run/suitex/work cylc-run/suitex/share/cycle cylc-run/suitex/share cylc-run/suitex; (cd ; rmdir -p cylc-run/suitex/work cylc-run/suitex/share/cycle cylc-run/suitex/share cylc-run 2>/dev/null || true)'

# rose suite-run: test for shared filesystem and create share/cycle directories
ssh -oBatchMode=yes -oConnectTimeout=10 -n hobo@wizard env ROSE_VERSION=2018.02.0 CYLC_VERSION=7.6.x bash -l -c '"$0" "$@"' rose suite-run -vv -n suitex --run=run --remote=uuid=231cd6a1-6d61-476d-96e1-4325ef9216fc,now-str=20180416T042319Z

# rose suite-run: install suite source directory to job host
rsync -a --exclude=.* --timeout=1800 --rsh=ssh -oBatchMode=yes -oConnectTimeout=10 --exclude=231cd6a1-6d61-476d-96e1-4325ef9216fc --exclude=log/231cd6a1-6d61-476d-96e1-4325ef9216fc --exclude=share/231cd6a1-6d61-476d-96e1-4325ef9216fc --exclude=share/cycle/231cd6a1-6d61-476d-96e1-4325ef9216fc --exclude=work/231cd6a1-6d61-476d-96e1-4325ef9216fc --exclude=/.* --exclude=/cylc-suite.db --exclude=/log --exclude=/log.* --exclude=/state --exclude=/share --exclude=/work ./ hobo@wizard:cylc-run/suitex
   # (internal rsync)
   ssh -oBatchMode=yes -oConnectTimeout=10 -l hobo wizard rsync --server -logDtpre.iLsfx --timeout=1800 . cylc-run/suitex
   # (internal rsync, back from hobo@wizard)
   rsync --server -logDtpre.iLsfx --timeout=1800 . cylc-run/suitex
\end{lstlisting}

\vspace{5mm}
Result:
\lstset{language=sh}
{\scriptsize
\dirtree{%
.1 /home/hjoliver/cylc-run/suitex.
.2 log->log.20180418T025047Z\DTcomment{\textbf{LOG DIRECTORIES}}.
.2 log.20180418T025047Z\DTcomment{log directory for current suite run}.
.2 suiter.rc.
.2 xxx\DTcomment{(any suite source sub-dirs or file)}.
.2 work\DTcomment{\textbf{JOB WORK DIRECTORIES}}.
.2 share\DTcomment{\textbf{SUITE SHARE DIRECTORY}}.
.3 cycle.
}
}

\subsection{server installs service directory}

\begin{itemize}
  \item server address and credentials, so that clients such as
    \lstinline=cylc message=  executed by jobs can connect
  \item done just before the first job is submitted to a remote, and at
    suite restart for the remotes of jobs running when the suite went
    down (server host, port, etc.\ may change at restart)
  \item uses SSH to invoke \lstinline=cylc remote-init= on 
    job hosts. If the remote command does not find a server-side UUID file
    (which would indicate a shared filesystem) it reads a tar archive of
    the service directory from stdin, and unpacks it to install.
\end{itemize}

{\color{blue}{\bf SSH-free?}
\begin{itemize}
  \item service directory can be installed with \lstinline=rsync=
  \item detect shared filesystems by attempting to \lstinline=rsync= the UUID
    file back from the job host?
\end{itemize}
}

\lstset{language=jobhosts}

\vspace{5mm}
\begin{lstlisting}
# cylc remote-init: install suite service directory
ssh -oBatchMode=yes -oConnectTimeout=10 hobo@wizard env CYLC_VERSION=7.6.x /opt/bin/cylc remote-init '066592b1-4525-48b5-b86e-da06eb2380d9' '$HOME/cylc-run/suitex'
\end{lstlisting}

Result:
{\scriptsize
\dirtree{%
.1 /home/hjoliver/cylc-run/suitex.
.2 .service\DTcomment{\textbf{SUITE SERVICE DIRECTORY}}. 
.3 contact\DTcomment{{\color{blue} server address information}}.
.3 passphrase\DTcomment{{\color{blue} suite passphrase}}.
.3 ssl.cert\DTcomment{{\color{blue} suite SSL certificate}}.
.2 log->log.20180418T025047Z\DTcomment{\textbf{LOG DIRECTORIES}}.
.2 log.20180418T025047Z\DTcomment{log directory for current suite run}.
.2 suiter.rc.
.2 xxx\DTcomment{(any suite source sub-dirs or file)}.
.2 work\DTcomment{\textbf{JOB WORK DIRECTORIES}}.
.2 share\DTcomment{\textbf{SUITE SHARE DIRECTORY}}.
.3 cycle.
}
}

\subsection{server submits jobs}
\begin{itemize}
  \item done when tasks are ready to run, for multiple jobs at once
  \item uses SSH to invoke \lstinline=cylc jobs-submit= on the
    remote - to read job scripts from stdin, write them to disk, and submit
    them to run
\end{itemize}

{\color{blue}{\bf SSH-free?}
\begin{itemize}
  \item just submit jobs to the local batch system
  \item with job host specified in batch system directives\footnote{Or for
    Slurm, a \lstinline=--cluster= option that needs to be used in subsequent
    poll and kill commands too.}
\end{itemize}
}

\lstset{language=jobhosts}

\vspace{5mm}
\begin{lstlisting}
# cylc jobs-submit: submit two jobs
ssh -oBatchMode=yes -oConnectTimeout=10 hobo@wizard env CYLC_VERSION=7.6.x /opt/bin/cylc jobs-submit '--remote-mode' '--' '$HOME/cylc-run/suitex/log/job' '1/remote-a/01' '1/remote-b/01'
\end{lstlisting}

Result:
{\scriptsize
\dirtree{%
.1 /home/hjoliver/cylc-run/suitex.
.2 .service\DTcomment{\textbf{SUITE SERVICE DIRECTORY}}. 
.3 contact\DTcomment{{\color{blue} server address information}}.
.3 passphrase\DTcomment{{\color{blue} suite passphrase}}.
.3 ssl.cert\DTcomment{{\color{blue} suite SSL certificate}}.
.2 log->log.20180418T025047Z\DTcomment{\textbf{LOG DIRECTORIES}}.
.2 log.20180418T025047Z\DTcomment{log directory for current suite run}.
.3 job\DTcomment{job logs (to be distinguished from \lstinline=log/suite/= on the suite host)}.
.4 1\DTcomment{cycle point}.
.5 remote-a\DTcomment{task name}.
.6 01\DTcomment{job submit number}.
.7 job\DTcomment{{\color{blue}job script}}.
.7 job.out\DTcomment{{\color{blue} job stdout}}.
.7 job.err\DTcomment{{\color{blue} job stderr}}.
.7 job.status\DTcomment{{\color{blue} job status}}.
.6 NN->0l\DTcomment{symlink to latest submit number}.
.5 remote-b\DTcomment{task name}.
.6 01\DTcomment{job submit number}.
.7 job\DTcomment{{\color{blue}job script}}.
.7 job.out\DTcomment{{\color{blue} job stdout}}.
.7 job.err\DTcomment{{\color{blue} job stderr}}.
.7 job.status\DTcomment{{\color{blue} job status}}.
.6 NN->0l\DTcomment{symlink to latest submit number}.
.2 suiter.rc.
.2 xxx\DTcomment{(any suite source sub-dirs or file)}.
.2 work\DTcomment{\textbf{JOB WORK DIRECTORIES}}.
.3 1\DTcomment{cycle point}.
.4 remote-a\DTcomment{task name}.
.5 xxx\DTcomment{(any files written by job to PWD)}.
.4 remote-b\DTcomment{task name}.
.5 xxx\DTcomment{(any files written by job to PWD)}.
.2 share\DTcomment{\textbf{SUITE SHARE DIRECTORY}}.
.3 cycle.
.3 xxx\DTcomment{(any job-created sub-dirs and files)}.
}
}

\subsection{server tracks job progress}

\begin{itemize}
  \item jobs send messages back to the server program on the suite host
    \begin{itemize}
      \item directly: client-server HTTPS over the network (requires service
        files installed - see above)
      \item indirectly: re-invoke clients on the suite host (requires reverse SSH)
    \end{itemize}
  \item OR server polls jobs at intervals (requires job polling - see below)
\end{itemize}

{\color{blue}{\bf SSH-free?}
\begin{itemize}
  \item for job messaging by HTTPS: no impact
  \item for server job polling: see polling description below
\end{itemize}
}

\subsection{user views job logs}

\begin{itemize}
  \item command \lstinline=cylc cat-log= via CLI or GUI, invokes itself over
    SSH to the remote
  \item suites will serve job logs in future, but this will still be needed
    (e.g.\ if the suite is down)
\end{itemize}

{\color{blue}{\bf SSH-free?}
\begin{itemize}
  \item \lstinline=rsync= the job logs back on demand, and view them locally
\end{itemize}
}

\vspace{5mm}
\begin{lstlisting}
# cylc cat-log: view a job log
ssh -oBatchMode=yes -oConnectTimeout=10 -n hobo@wizard env CYLC_VERSION=7.6.x /opt/bin/cylc cat-log --remote-arg='$HOME/cylc-run/suitex/log/job/1/remote-a/NN/job.out' --remote-arg=cat --remote-arg='tail -n +1 -F %(filename)s' suitex
\end{lstlisting}


\subsection{server cancels or kills jobs}

\begin{itemize}
  \item done automatically or via user command \lstinline=cylc kill=, for
    multiple jobs at once
  \item uses SSH to invoke \lstinline=cylc jobs-kill= on the
    remote, with job log paths on the command line. Reads job ID from the
    job status file.
\end{itemize}

{\color{blue}{\bf SSH-free?}
\begin{itemize}
  \item read job ID from the new job management file (see below) and kill via
    the local batch system
\end{itemize}
}

\vspace{5mm}
    \begin{lstlisting}
# cylc jobs-kill: kill two jobs
ssh -oBatchMode=yes -oConnectTimeout=10 hobo@wizard env CYLC_VERSION=7.6.x /opt/bin/cylc jobs-kill '--' '$HOME/cylc-run/suitex/log/job' '1/remote-a/01' '1/remote-b/01'
    \end{lstlisting}

\subsection{server polls jobs}

\begin{itemize}
  \item done automatically or via user command \lstinline=cylc poll=, for
    multiple jobs at once
  \item uses SSH to invoke \lstinline=cylc jobs-poll= on the
    remote, with job log paths on the command line. Reads job ID from the
    job status file.
\end{itemize}

{\color{blue}{\bf SSH-free?}
\begin{itemize}
  \item read job ID from the new job management file (see below)
  \item \lstinline=rsync= the new job message file (see below) from the job
    host to determine job exit status (if any)
  \item then if needed, query the local batch system
\end{itemize}
}

\vspace{5mm}
    \begin{lstlisting}
# cylc jobs-poll: poll two jobs
ssh -oBatchMode=yes -oConnectTimeout=10 hobo@wizard env CYLC_VERSION=7.6.x /opt/bin/cylc jobs-poll '--' '$HOME/cylc-run/suitex/log/job' '1/remote-a/01' '1/remote-b/01'
    \end{lstlisting}


\subsection{server retrieves jobs logs}

\begin{itemize}
  \item done at job completion, according to global config
  \item uses \lstinline=rsync=
\end{itemize}

{\color{blue}{\bf SSH-free?}
\begin{itemize}
  \item (already done via \lstinline=rsync=)
\end{itemize}
}

\vspace{5mm}
    \begin{lstlisting}
# rsync: retrieve two job logs
rsync -a --rsh=ssh -oBatchMode=yes -oConnectTimeout=10 --include=/1 --include=/1/remote-a --include=/1/remote-a/01 --include=/1/remote-a/01/** --include=/1/remote-b --include=/1/remote-b/01 --include=/1/remote-b/01/** --exclude=/** hobo@wizard:$HOME/cylc-run/suitex/log/job/ /home/vagrant/cylc-run/suitex/log/job/
   # (internal rsync)
   ssh -oBatchMode=yes -oConnectTimeout=10 -l hobo wizard rsync --server --sender -logDtpre.iLsfx . $HOME/cylc-run/suitex/log/job/
   # (internal rsync, back from hobo@wizard)
   rsync --server --sender -logDtpre.iLsfx . /home/hobo/cylc-run/suitex/log/job/
    \end{lstlisting}

\subsection{server tidies job remote at shutdown}

\begin{itemize}
  \item removes \lstinline=.service/contact= so that clients won't repeatedly
    try to connect
\end{itemize}

{\color{blue}{\bf SSH-free?}
\begin{itemize}
  \item remove with \lstinline=rsync='s delete option and a modified source
    directory on the suite host?
\end{itemize}
}

\vspace{5mm}
    \begin{lstlisting}
# cylc remote-tidy: remove the remote suite contact file
ssh -oBatchMode=yes -oConnectTimeout=10 hobo@wizard env CYLC_VERSION=7.6.x /opt/bin/cylc remote-tidy '$HOME/cylc-run/suitex'
    \end{lstlisting}

\section{SSH-free Job Management: Implementation Plan}
\label{Implementation}

This section lays out a path to implementation of SSH-free job management
with a local batch system on the suite host, in light of the blue
{\color{blue}SSH-free?} comments above.

\begin{itemize}
  \item suite directories must be created on the job host, and files installed
    there, using \lstinline=rsync=
    \begin{itemize}
      \item by \lstinline=rose/cylc suite-run= or at
        remote init time (before the first job submission there)
    \end{itemize}

  \item detect shared filesystems by attempting to \lstinline=rsync= the server
    UUID file back from the job host

  \item \lstinline=cylc jobs-submit=, \lstinline=cylc jobs-poll=, and
    \lstinline=cylc jobs-kill= must run on the {\em batch system host}
    \begin{itemize}
      \item in SSH mode this is the job host
      \item in SSH-free mode this is the suite host
      \item cylc now needs to distinguish batch system host as well as suite
        host and job host
    \end{itemize}

  \item the submit command also writes to the {\bf job status file}, and the
    poll and kill commands also read from it, on the job host. In SSH-free mode
    the status files are on the wrong host. We can get around this problem by
    splitting the existing job status files into two new job files:
    \begin{itemize}
      \item {\bf job message files} written by jobs, on the job host, for run
        status and other messages sent back to the server, and job exit status
      \item {\bf job management files} written by job management commands on
        the batch system host (for batch system name and job ID)
    \end{itemize}

  \item job management commands can write to and read from the new job
    management files as needed (for batch system name and job ID) because they
    always run on the right host (in both modes).

  \item the job poll command still needs to read job exit status (if any) from
    the job message file on the job host, before deciding whether or not to
    query the batch system
    \begin{itemize}
      \item for SSH-based job management this is trivial (same host)
      \item in SSH-free mode the poll command must \lstinline=rsync=
        the job message file back to the suite host
    \end{itemize}

  \item if possible we should allow combined mode operation: the suite server
    program should use ssh to interact with a batch system host that is not the
    job host

  \item finally, how to configure SSH-free operation: can it be determined
    automatically (when specified job and batch system hosts differ) or does it
    need an explicit global config option?
\end{itemize}

\section{Appendix: Other Use of SSH in Cylc}

Two other uses of SSH in cylc do not impact on the topic of this document:

\begin{itemize}
  \item see if a suite is running on another host with a shared
    filesystem - see \lstinline=detect_old_contact_file()= in
    \lstinline=lib/cylc/suite_srv_files_mgr.py=
  \item cat content of a remote service file over SSH, if possible, for clients
    on that do not have suite credentials installed - see
    \lstinline=_load_remote_item()= in \lstinline=suite_srv_files_mgr.py=
\end{itemize}

\end{document}
